\documentclass[11pt]{article}

\usepackage[T1]{fontenc}
\usepackage[utf8]{inputenc}
\usepackage{lmodern}
\usepackage[margin=1.04in]{geometry}
\usepackage{amsmath,amssymb,mathtools,amsthm}
\usepackage{aliascnt}
\usepackage{mathrsfs}
\usepackage{microtype}
\usepackage{enumitem}
\usepackage{xcolor}
\usepackage{booktabs}
\usepackage{url}
\usepackage{hyperref}
\usepackage[capitalize,noabbrev]{cleveref}

\hypersetup{
  colorlinks=true,
  linkcolor=blue!45!black,
  citecolor=green!35!black,
  urlcolor=blue!55!black,
  pdftitle={Supplementary Results on Compact Exponential Periods: Quadratic Normal Forms, Tensor-Square Residues, and Entire Exponential Integrals},
  pdfauthor={Christopher D. Long; Antoine-Auguste Le Blanc},
  pdfsubject={Quadratic exponential integrals, tensor-square residues, entire exponential integrals},
  pdfkeywords={E-function, algebraic independence, exponential period, tensor square, quartic polynomial, entire exponential integral}
}

\numberwithin{equation}{section}
\setlength{\emergencystretch}{3em}

\newtheorem{theorem}{Theorem}[section]
\newaliascnt{proposition}{theorem}
\newtheorem{proposition}[proposition]{Proposition}
\aliascntresetthe{proposition}
\newaliascnt{lemma}{theorem}
\newtheorem{lemma}[lemma]{Lemma}
\aliascntresetthe{lemma}
\newaliascnt{corollary}{theorem}
\newtheorem{corollary}[corollary]{Corollary}
\aliascntresetthe{corollary}
\theoremstyle{definition}
\newaliascnt{definition}{theorem}
\newtheorem{definition}[definition]{Definition}
\aliascntresetthe{definition}
\theoremstyle{remark}
\newaliascnt{remark}{theorem}
\newtheorem{remark}[remark]{Remark}
\aliascntresetthe{remark}

\newcommand{\C}{\mathbb C}
\newcommand{\Qbar}{\overline{\mathbb Q}}
\newcommand{\Aone}{\mathbb A^1}
\newcommand{\Spec}{\operatorname{Spec}}
\newcommand{\Span}{\operatorname{span}}
\newcommand{\Sym}{\operatorname{Sym}}
\newcommand{\End}{\operatorname{End}}
\newcommand{\ord}{\operatorname{ord}}
\newcommand{\Ein}{\operatorname{Ein}}
\newcommand{\Acal}{\mathcal A}
\newcommand{\Mcal}{\mathcal M}
\newcommand{\dd}{\,d}

\title{Supplementary Results on Compact Exponential Periods:\\
Quadratic Normal Forms, Tensor-Square Residues, and Entire Exponential Integrals}
\author{Christopher D. Long \and Antoine-Auguste Le Blanc\thanks{Le Blanc cryptographic author identifier: \texttt{b6048125\allowbreak{}30b3535b\allowbreak{}23d6dbb0\allowbreak{}f5abe32d\allowbreak{}1ed1a9de\allowbreak{}35fb4f7d\allowbreak{}04d62cff\allowbreak{}08ea19c8}}}
\date{August 2026}

\begin{document}
\maketitle

\begin{center}
\small
\texttt{galizur@gmail.com}
\end{center}

\begin{abstract}
This note preserves three independent developments omitted from the streamlined
paper \emph{A Polynomial Kontsevich--Zagier Theorem for Compact Exponential
Periods on the Affine Line}.  First, for
\[
  H_\beta=\int_0^1e^{\beta x^2}\,dx,
  \qquad \beta\in\overline{\mathbb Q}^{\times},
\]
we prove algebraic independence of all finite families over the field
generated by algebraic exponential values and derive an explicit polynomial
presentation of the formal period algebra when the polynomial in the
exponent has degree at most two.  Second, we analyze rational horizontal
endomorphisms of tensor squares of polynomial moment connections.  Constancy
holds through cubic polynomials; for quartics the first integral residue
resonance occurs exactly for affine pure quartics, and its endomorphism
algebra is computed explicitly.  Third, we prove algebraic independence of
finite families of values of the entire exponential integral
$\operatorname{Ein}$ over the same exponential field.  A conjugation and
field-intersection argument implies that $\pi$ and $\operatorname{Ein}(1+i)$
are algebraically independent.
\end{abstract}

\medskip
\noindent\textit{2020 Mathematics Subject Classification.}
Primary 11J85; Secondary 11J81, 11J91, 12H05, 34M35.

\noindent\textit{Keywords.}
Algebraic independence, $E$-functions, exponential periods, quadratic
exponential integrals, tensor-square residues, entire exponential integral.

\section{Purpose and notation}

The principal Paper III is most transparent when it is devoted solely to the
all-degree formal-period theorem \cite{LongPolynomialKZ}.  The results below
are mathematically useful but logically independent of that proof.  They are
collected here so that they can be audited, cited, or separated into
standalone notes without enlarging the main dependency chain.

Put $k=\Qbar$ and
\begin{equation}\label{eq:Efield}
  \mathbb E=k(e^\lambda:\lambda\in k),
\end{equation}
the field generated by all exponential values at algebraic arguments.  Any
single algebraic relation over $\mathbb E$ involves only finitely many such
exponentials, so every proof reduces to a finitely generated additive
subgroup of $k$.

\section{Differential-algebra preliminaries}

\subsection{\texorpdfstring{$E$}{E}-functions and specialization}

We use Beukers' refinement of the Siegel--Shidlovskii theorem
\cite{Beukers}.  In particular, at every nonsingular nonzero algebraic point,
the transcendence degree of the values of an $E$-function vector equals the
transcendence degree of the corresponding functions over $k(z)$.

\begin{theorem}[Lindemann--Weierstrass, linear form]
\label{thm:LW}
If $\lambda_1,\ldots,\lambda_s\in k$ are distinct, then
$e^{\lambda_1},\ldots,e^{\lambda_s}$ are linearly independent over $k$.
Consequently, if $\Gamma\subset k$ is a finitely generated additive subgroup
of rank $r$, then
\[
  \operatorname{trdeg}_k k(e^\lambda:\lambda\in\Gamma)=r.
\]
\end{theorem}

\begin{proof}
The linear assertion is classical; see \cite[Chapter~1]{Baker}.  Choose a
$\mathbb Z$-basis $\omega_1,\ldots,\omega_r$ of $\Gamma$.  A polynomial
relation among the $e^{\omega_j}$ expands into a linear relation among
exponentials of distinct algebraic numbers.
\end{proof}

\subsection{Exponential fields and character extraction}

Let $u$ be a complex variable, let $d\in\{1,2\}$, and let
$\Gamma\subset\C$ be a finitely generated additive subgroup with basis
$\omega_1,\ldots,\omega_r$.  Put
\begin{equation}\label{eq:K-Gamma-d}
  K_{\Gamma,d}
  =\C\bigl(u,e^{\omega_1u^d},\ldots,e^{\omega_ru^d}\bigr).
\end{equation}
Distinct polynomial exponentials are linearly independent over $\C(u)$, so
this is a rational function field in the displayed exponential generators.
Its constants for $d/du$ are $\C$.  The torus
$\mathbf T=(\C^\times)^r$ acts differentially by independently rescaling the
exponential generators and fixes $\C(u)$.

\begin{lemma}[Extraction of one exponential character]
\label{lem:character-extraction}
Suppose $g\in K_{\Gamma,d}$ and
\[
  g'=\sum_{\mu\in S}h_\mu(u)e^{\mu u^d},
  \qquad h_\mu(u)\in\C(u),
\]
where $S\subset\Gamma$ is finite.  For every nonzero $\mu\in S$ with
$h_\mu\ne0$, there is $r_\mu(u)\in\C(u)$ such that
\[
  \frac{d}{du}\left(r_\mu(u)e^{\mu u^d}\right)
  =h_\mu(u)e^{\mu u^d}.
\]
\end{lemma}

\begin{proof}
Work in $K_{\Gamma,d}/\C$.  Differentiation is injective and torus
equivariant.  The derivative of the torus orbit of the class of $g$ lies in
the finite-dimensional span of the displayed characters; hence the orbit
span of the class of $g$ is finite-dimensional.  Decompose it into torus
weight spaces.  A representative of the nonzero $\mu$-weight can be chosen
to be an exact semi-invariant because a torus is linearly reductive.  Dividing
that representative by $e^{\mu u^d}$ gives a torus-invariant element, hence an
element of $\C(u)$.
\end{proof}

\begin{theorem}[Kolchin--Ostrowski]\label{thm:KO}
Let $K$ be a differential field of characteristic zero with algebraically
closed constant field $C$.  Suppose $y_1,\ldots,y_m$ lie in a differential
extension with the same constants and satisfy $y_j'\in K$.  They are
algebraically dependent over $K$ if and only if a nonzero constant linear
combination of them belongs to $K$.
\end{theorem}

This is the antiderivative theorem of Kolchin and Ostrowski; see
\cite{Kolchin,vanDerPutSinger}.

\begin{lemma}[Two rational differential equations]
\label{lem:rational-odes}
Let $\alpha,\beta\in\C^\times$ and $c\in\C^\times$.
\begin{enumerate}[label=\textup{(\roman*)},leftmargin=2.4em]
\item The equation
\[
  r'(u)+2\beta u r(u)=c
\]
has no solution $r\in\C(u)$.
\item The equation
\[
  r'(u)-\alpha r(u)=\frac{c}{u}
\]
has no solution $r\in\C(u)$.
\end{enumerate}
\end{lemma}

\begin{proof}
In (i), pole-order comparison excludes every finite pole, so $r$ is a
polynomial.  The term $2\beta ur$ has degree $\deg r+1$, whereas $r'$ has
degree at most $\deg r-1$, so the sum cannot be constant.  In (ii), finite poles away from zero are
excluded similarly.  A pole at zero makes $r'$ have strictly larger order
than both $\alpha r$ and $c/u$; regularity at zero is also impossible because
the right side has a pole.
\end{proof}

\section{Quadratic exponential integral values}
\label{sec:quadratic}

For $\beta\in\C^\times$, define
\begin{equation}\label{eq:Hbeta}
  H_\beta(z)=\int_0^1e^{\beta z x^2}\,dx,
  \qquad H_\beta=H_\beta(1).
\end{equation}
Termwise integration and one integration by parts give
\begin{equation}\label{eq:H-series-ode}
  H_\beta(z)=\sum_{n\ge0}\frac{\beta^n}{2n+1}\frac{z^n}{n!},
  \qquad
  2zH_\beta'(z)+H_\beta(z)=e^{\beta z}.
\end{equation}
The standard common-denominator estimate for $(2n+1)^{-1}\beta^n$ shows that
$H_\beta$ is an $E$-function when $\beta$ is algebraic.

\begin{theorem}[Functional independence of the quadratic integrals]
\label{thm:H-functional}
For distinct $\beta_1,\ldots,\beta_m\in\C^\times$, the functions
$H_{\beta_1}(z),\ldots,H_{\beta_m}(z)$ are algebraically independent over
the field
\[
  \C\bigl(z,e^{\lambda z}:\lambda\in\C\bigr).
\]
\end{theorem}

\begin{proof}
Any proposed algebraic relation uses only finitely many exponential
functions.  Let $\Gamma$ be the additive subgroup generated by their indices
and by the $\beta_j$.  Put $z=w^2$ and
\[
  Y_\beta(w)=wH_\beta(w^2)=\int_0^w e^{\beta u^2}\,du.
\]
Then $Y_\beta'=e^{\beta w^2}$, and the pullback of the coefficient field is
contained in $K_{\Gamma,2}$.

If the $Y_{\beta_j}$ were algebraically dependent, \cref{thm:KO} would give
constants $c_j$, not all zero, and $g\in K_{\Gamma,2}$ with
\[
  g'=\sum_jc_je^{\beta_jw^2}.
\]
For every $j$ with $c_j\ne0$, \cref{lem:character-extraction} produces
$r_j\in\C(w)$ satisfying
\[
  r_j'+2\beta_jwr_j=c_j,
\]
contrary to \cref{lem:rational-odes}(i).
\end{proof}

\begin{theorem}[Quadratic value independence]
\label{thm:H-values}
For distinct $\beta_1,\ldots,\beta_m\in k^\times$, the values
$H_{\beta_1},\ldots,H_{\beta_m}$ are algebraically independent over
$\mathbb E$.
\end{theorem}

\begin{proof}
A proposed algebraic relation over $\mathbb E$ involves only finitely many
exponential values.  Let $\Gamma\subset k$ be generated by their indices and
the $\beta_j$.  Form the $E$-function vector consisting of the $H_{\beta_j}$
and a basis of exponential functions for $\Gamma$.  By
\cref{thm:H-functional}, its functional transcendence degree is
$\operatorname{rank}\Gamma+m$.  Beukers' theorem gives the same
transcendence degree at $z=1$, while \cref{thm:LW} gives transcendence degree
$\operatorname{rank}\Gamma$ for the exponential-value subfield.  The $m$
quadratic values are therefore algebraically independent over that subfield.
\end{proof}

Let $\mathscr E_k=k[(k,+)]$ be the group algebra with basis
$\varepsilon_\lambda$.  Evaluation sends $\varepsilon_\lambda$ to
$e^\lambda$.

\begin{corollary}[Universal quadratic value ring]
\label{cor:universal-H-ring}
Evaluation induces an injective homomorphism
\begin{equation}\label{eq:universal-H-ring}
  \mathscr E_k[T_\beta:\beta\in k^\times]
  \longrightarrow\C,
  \qquad
  \varepsilon_\lambda\longmapsto e^\lambda,
  \quad
  T_\beta\longmapsto H_\beta.
\end{equation}
\end{corollary}

\begin{proof}
Every polynomial involves finitely many $T_\beta$ and finitely many group
indices.  Apply \cref{thm:H-values} and Lindemann--Weierstrass.
\end{proof}

\subsection{The degree-at-most-two formal algebra}

Let $\Acal_{\Aone,\le2}^{\mathrm{KZ,cpt}}$ be the subalgebra of the formal
algebra in \cite{LongPolynomialKZ} generated by symbols
$\mathsf P(q,r,\Delta)$ with $\deg q\le2$ and by the symbols
$\mathsf E(\lambda)$.  For $\beta\in k^\times$, put
\[
  \mathsf H_\beta
  =\mathsf P(\beta x^2,1,[1]-[0]).
\]

\begin{lemma}[Quadratic normal form]\label{lem:quadratic-normal-form}
Let
\[
  q(x)=a(x-h)^2+c,
  \qquad a\in k^\times,
  \quad h,c\in k,
\]
and let $\Delta=\sum_\alpha d_\alpha[\alpha]$.  Then
\begin{equation}\label{eq:quadratic-normal-form}
  \mathsf P(q,1,\Delta)
  =\mathsf E(c)
   \sum_\alpha d_\alpha(\alpha-h)
      \mathsf H_{a(\alpha-h)^2},
\end{equation}
where the term with $\alpha=h$ is zero.
\end{lemma}

\begin{proof}
Because $\sum_\alpha d_\alpha=0$,
$\Delta=\sum_\alpha d_\alpha([\alpha]-[h])$.  On the path from $h$ to
$\alpha$, substitute $x=h+(\alpha-h)t$, and then use the constant-shift
identity to factor $\mathsf E(c)$.
\end{proof}

\begin{theorem}[Explicit degree-two period algebra]
\label{thm:degree-two-algebra}
The homomorphism
\[
  \mathscr E_k[T_\beta:\beta\in k^\times]
  \longrightarrow
  \Acal_{\Aone,\le2}^{\mathrm{KZ,cpt}},
  \qquad
  \varepsilon_\lambda\longmapsto\mathsf E(\lambda),
  \quad
  T_\beta\longmapsto\mathsf H_\beta,
\]
is an isomorphism, and the numerical period map on
$\Acal_{\Aone,\le2}^{\mathrm{KZ,cpt}}$ is injective.
\end{theorem}

\begin{proof}
The twisted reduction of \cite{LongLinearKZ} reduces a polynomial
differential for a quadratic $q$ to a constant multiple of $dx$ plus
zero-dimensional terms.  Constant and linear polynomials in the exponent
also reduce to zero-dimensional terms.  Completing the square and applying
\cref{lem:quadratic-normal-form} proves surjectivity.  The composite with the
numerical period map is \eqref{eq:universal-H-ring}, which is injective by
\cref{cor:universal-H-ring}; hence the displayed homomorphism and the period
map are injective.
\end{proof}

\section{Tensor-square residue analysis}
\label{sec:tensor}

Let $\Mcal_q$ be the homogeneous moment module associated in
\cite{LongEIT} with a polynomial $q$ of degree $n\ge2$.  In a moment basis its
connection matrix is
\[
  A=C+\frac Dz,
\]
where $C$ is semisimple,
\[
  \Spec(D)=\left\{-\frac1n,\ldots,-\frac{n-1}{n}\right\},
\]
and the compressed residue at a critical point of multiplicity $m$ has
spectrum
\[
  -\frac m{m+1},-\frac{m-1}{m+1},\ldots,-\frac1{m+1}.
\]
On a tensor square,
\[
  C^{(2)}=C\otimes I+I\otimes C,
  \qquad
  D^{(2)}=D\otimes I+I\otimes D.
\]

\begin{lemma}[Compressed tensor residues]
\label{lem:tensor-compression}
For an eigenvalue $\Lambda$ of $C^{(2)}$,
\[
  \ker(C^{(2)}-\Lambda I)
  =\bigoplus_{\lambda+\mu=\Lambda}V_\lambda\otimes V_\mu,
\]
and the compression of $D^{(2)}$ to this eigenspace is
\[
  \bigoplus_{\lambda+\mu=\Lambda}
  \left(\overline D_\lambda\otimes I
        +I\otimes\overline D_\mu\right).
\]
\end{lemma}

\begin{proof}
Use the spectral projectors of the semisimple matrix $C$.  Compression of
$D\otimes I$ retains only the $V_\lambda$ component of $D$ in the first
factor, and similarly for the second factor.  Distinct ordered pairs do not
mix, even when their eigenvalue sums coincide.
\end{proof}

\begin{theorem}[Tensor-square rigidity through degree three]
\label{thm:degree23}
Let $M$ be a finite direct sum of moment modules attached to polynomials of
degree two or three.  Every rational horizontal endomorphism of
$M^{\otimes2}$ and of $\Sym^2M$ is constant in the tensor moment basis.
\end{theorem}

\begin{proof}
The possible residue eigenvalues of $M$ are
$-1/2,-1/3,-2/3$.  Those of $D^{(2)}$ lie in
$[-4/3,-2/3]$, an interval of width $2/3$.  A pole at zero of order
$m\ge1$ would make $-m$ an eigenvalue of $\operatorname{ad}D^{(2)}$, which
is impossible; poles away from zero are excluded by order comparison.
Thus a rational horizontal endomorphism is polynomial.

If its degree at infinity were $d\ge1$, compression of the two leading
coefficient equations to an eigenspace of $C^{(2)}$ would make $d$ an
eigenvalue of the commutator with the compressed tensor residue.  By
\cref{lem:tensor-compression}, that residue again has spectrum in an
interval of width $2/3$.  Hence $d$ cannot be a positive integer.  The
symmetric square is a horizontal direct summand of the tensor square.
\end{proof}

\subsection{The first quartic resonance}

After an algebraic affine change of the source, a quartic has the form
\begin{equation}\label{eq:quartic}
  q(x)=\frac{x^4}{4}+ax^2+bx+c.
\end{equation}
It is affine-equivalent to a pure quartic
$\alpha(x-h)^4+\beta$ exactly when $a=b=0$ in this normalization.

We use the row-module convention, so a horizontal endomorphism $Q$ satisfies
\[
  Q'=[Q,C+D/z].
\]
In the basis
\[
  v_0=1,
  \qquad v_1=x,
  \qquad v_2=x^2+a,
\]
the row-connection matrices are
\begin{equation}\label{eq:quartic-CD}
  D=\operatorname{diag}\left(-\frac14,-\frac12,-\frac34\right),
\end{equation}
and
\[
  C=
  \begin{pmatrix}
  c-\frac{a^2}{2}&-\frac{5ab}{4}&\frac{a^3}{2}-\frac{3b^2}{4}\\[1mm]
  \frac{3b}{4}&c-a^2&-\frac{5ab}{4}\\[1mm]
  \frac a2&\frac{3b}{4}&c-\frac{a^2}{2}
  \end{pmatrix}.
\]
These formulas follow by dividing $qv_j$ by $q'$ and taking the remainder.
On $\Sym^2V$, in the ordered product basis
\[
  v_0^2,
  \ v_0v_1,
  \ v_0v_2,
  \ v_1^2,
  \ v_1v_2,
  \ v_2^2,
\]
one has
\begin{equation}\label{eq:quartic-D2}
  D^{(2)}=
  \operatorname{diag}\left(-\frac12,-\frac34,-1,-1,-\frac54,-\frac32\right).
\end{equation}

\begin{lemma}[Pole exclusion for a non-pure quartic]
\label{lem:quartic-pole}
If $(a,b)\ne(0,0)$, a rational horizontal endomorphism of the tensor square
has no pole at zero.
\end{lemma}

\begin{proof}
The width of the spectrum in \eqref{eq:quartic-D2} is one.  A pole can only
have order one and, after projection to the symmetric square, its leading
term is the map $E_{0,5}$ from the lowest residue line to the highest.  Write
\[
  Q=z^{-1}E_{0,5}+Q_0.
\]
The coefficient equations are
\[
  [Q_0,D^{(2)}]=-[E_{0,5},C^{(2)}],
  \qquad
  [Q_0,C^{(2)}]=0.
\]
Solving the first equation entrywise gives
\[
Q_0=
\begin{pmatrix}
 u_0&0&a&0&b&0\\
 0&u_1&0&0&0&-2b\\
 0&0&u_2&u_3&0&-2a\\
 0&0&u_4&u_5&0&0\\
 0&0&0&0&u_6&0\\
 0&0&0&0&0&u_7
\end{pmatrix}.
\]
The $(0,0)$ entry of the second equation is $a^2$, so $a=0$.  If
$b\ne0$, four further entries give
\[
  u_1=u_0,
  \qquad u_0-u_2=1,
  \qquad u_3=-2,
  \qquad u_1-u_2-u_3=0,
\]
a contradiction.  Hence $b=0$ as well.
\end{proof}

\begin{theorem}[Quartic tensor-square classification]
\label{thm:quartic-classification}
Let $q$ be a quartic polynomial.
\begin{enumerate}[label=\textup{(\roman*)},leftmargin=2.4em]
\item If $q$ is not affine-equivalent to a pure quartic, every rational
horizontal endomorphism of $\Mcal_q^{\otimes2}$ is constant.
\item If $q$ is pure, then
\begin{align*}
  \End^\nabla(\Mcal_q^{\otimes2})
  &\simeq M_2(\C)\oplus M_2(\C)\oplus M_3(\C)\oplus M_2(\C),\\
  \End^\nabla(\Sym^2\Mcal_q)
  &\simeq M_2(\C)\oplus\C\oplus M_2(\C)\oplus\C.
\end{align*}
The nonconstant entries are the factors $z$ and $z^{-1}$ joining the two
extreme residue lines.
\end{enumerate}
\end{theorem}

\begin{proof}
If $q$ is not pure, its critical multiplicities are $1+1+1$ or $2+1$.
The compressed tensor residues then lie in an interval of width less than
one, so positive degree at infinity is excluded exactly as in
\cref{thm:degree23}; \cref{lem:quartic-pole} excludes the only possible pole
at zero.

For a pure quartic, $C$ is scalar.  The tensor residue weights are
\[
\begin{array}{c|ccccc}
\text{weight}&-\frac12&-\frac34&-1&-\frac54&-\frac32\\ \hline
\text{multiplicity}&1&2&3&2&1.
\end{array}
\]
An entry joining weights $d_i,d_j$ is a constant multiple of
$z^{d_j-d_i}$ and is rational exactly when the difference is integral.
Grouping weights modulo $\mathbb Z$ gives block dimensions $2,2,3,2$ on the
full tensor square.  On the symmetric square the multiplicities are
$1,1,2,1,1$, giving block dimensions $2,1,2,1$.
\end{proof}

\begin{remark}
The quartic endomorphisms are symmetries of a homogeneous tensor connection,
not polynomial equations for the inhomogeneous moment orbit.  The affine
linearity theorem in \cite{LongPolynomialKZ} proves that they create no new
polynomial relations among period values.
\end{remark}

\section{Entire exponential-integral values}
\label{sec:Ein}

Define
\begin{equation}\label{eq:Ein}
  \Ein(z)=\int_0^z\frac{1-e^{-t}}{t}\,dt
  =\sum_{n\ge1}\frac{(-1)^{n-1}z^n}{n\,n!}.
\end{equation}
For $\alpha\in\C^\times$, put $E_\alpha(z)=\Ein(\alpha z)$.  Then
\begin{equation}\label{eq:Ein-ode}
  zE_\alpha'(z)=1-e^{-\alpha z}.
\end{equation}
When $\alpha$ is algebraic, the usual denominator estimate for
$\alpha^n/n$ and \eqref{eq:Ein-ode} show that $E_\alpha$ is an $E$-function.

\begin{theorem}[Functional independence of entire exponential integrals]
\label{thm:Ein-functional}
For distinct $\alpha_1,\ldots,\alpha_m\in\C^\times$, the functions
$E_{\alpha_1}(z),\ldots,E_{\alpha_m}(z)$ are algebraically independent over
\[
  \C\bigl(z,e^{\lambda z}:\lambda\in\C\bigr).
\]
\end{theorem}

\begin{proof}
Any proposed relation uses only finitely many exponentials; let $\Gamma$ be
the subgroup generated by their indices and the $\alpha_j$.  If the displayed
functions were algebraically dependent over $K_{\Gamma,1}$,
\cref{thm:KO} would give constants $c_j$, not all zero, and
$g\in K_{\Gamma,1}$ such that
\[
  g'=\frac{\sum_jc_j}{z}
     -\sum_j\frac{c_j}{z}e^{-\alpha_jz}.
\]
For any $j$ with $c_j\ne0$, character extraction gives
$r_j\in\C(z)$ satisfying
\[
  r_j'-\alpha_jr_j=-\frac{c_j}{z},
\]
contrary to \cref{lem:rational-odes}(ii).
\end{proof}

\begin{theorem}[Algebraic independence of $\Ein$-values]
\label{thm:Ein-values}
For distinct $\alpha_1,\ldots,\alpha_m\in k^\times$, the values
\[
  \Ein(\alpha_1),\ldots,\Ein(\alpha_m)
\]
are algebraically independent over $\mathbb E$.
\end{theorem}

\begin{proof}
A relation over $\mathbb E$ involves finitely many exponential values.  Form
the $E$-function vector consisting of the $E_{\alpha_j}$, the constant
function $1$, and exponential functions whose indices generate the finite
coefficient field.  Apply \cref{thm:Ein-functional}, Beukers' transcendence-
degree equality at $z=1$, and \cref{thm:LW}.
\end{proof}

The result is compatible with the classical first-order theory of
$E$-functions \cite{Salikhov} and with the structural work of Rivoal and
Roques \cite{RivoalRoques}.  The effective computation of relation ideals for
a fixed tuple is treated by Fischler and Rivoal
\cite{FischlerRivoalEffective}.

\section{Conjugation and an ordinary period}

\begin{lemma}[Intersection of one-variable algebraic closures]
\label{lem:field-intersection}
Let $k$ be algebraically closed, let $x,y$ be algebraically independent over
$k$, and let $t$ lie in a common overfield.  If $t$ is algebraic over both
$k(x)$ and $k(y)$, then $t\in k$.
\end{lemma}

\begin{proof}
If $t\notin k$, then $t$ is transcendental over $k$.  Algebraic dependence
makes both $x$ and $y$ algebraic over $k(t)$, so $k(x,y,t)$ has
transcendence degree one.  This contradicts the inclusion of $k(x,y)$, which
has transcendence degree two.
\end{proof}

\begin{theorem}[Conjugate-pair principle]\label{thm:conjugate-pair}
Let $F(z)=\sum_{n\ge0}a_nz^n$ have real algebraic coefficients and converge
at $\alpha$ and $\overline\alpha$.  If
$F(\alpha),F(\overline\alpha)$ are algebraically independent over $k$, then
every real transcendental number $\tau$ is algebraically independent from
$F(\alpha)$ over $k$.
\end{theorem}

\begin{proof}
If $\tau$ and $F(\alpha)$ were algebraically dependent, then $\tau$ would be
algebraic over $k(F(\alpha))$.  Complex conjugation fixes $\tau$, preserves
$k$, and sends $F(\alpha)$ to $F(\overline\alpha)$, so $\tau$ would also be
algebraic over $k(F(\overline\alpha))$.  Apply
\cref{lem:field-intersection}.
\end{proof}

\begin{theorem}[Algebraic independence of $\pi$ and $\Ein(1+i)$]
\label{thm:pi-Ein}
The numbers $\pi$ and $\Ein(1+i)$ are algebraically independent over
$\Qbar$.
\end{theorem}

\begin{proof}
By \cref{thm:Ein-values}, the conjugate pair
$\Ein(1+i),\Ein(1-i)$ is algebraically independent over $k$.  The Taylor
coefficients of $\Ein$ are real, and $\pi$ is real and transcendental.  Apply
\cref{thm:conjugate-pair}.
\end{proof}

The second number is itself a compact two-dimensional exponential period:
\begin{equation}\label{eq:Ein-square}
  \Ein(\alpha)
  =\alpha\int_0^1\int_0^1e^{-\alpha xy}\,dy\,dx.
\end{equation}
Thus \cref{thm:pi-Ein} is a mixed algebraic-independence result between an
ordinary period and a compact exponential period.  It does not assert a
formal-period theorem for the category generated by all ordinary and
exponential periods.

\section*{Acknowledgments and AI-assisted research disclosure}

This note preserves material developed alongside \cite{LongEIT,LongLinearKZ,
LongPolynomialKZ}.  During its preparation the author used OpenAI ChatGPT
5.6 Pro and Anthropic Claude Opus 5 for proof exploration, adversarial
checking, symbolic verification of the quartic matrices, bibliographic
verification, and language editing.  The AI systems are not authors.  The
authors reviewed the arguments and bear full responsibility for every
statement, proof, citation, and remaining error.

\small
\begin{thebibliography}{99}

\bibitem{Baker}
A.~Baker,
\emph{Transcendental Number Theory},
Cambridge Mathematical Library,
Cambridge University Press, Cambridge, 1990.

\bibitem{Beukers}
F.~Beukers,
\emph{A refined version of the Siegel--Shidlovskii theorem},
Ann. of Math. (2) \textbf{163} (2006), no.~1, 369--379.
\href{https://doi.org/10.4007/annals.2006.163.369}{doi:10.4007/annals.2006.163.369}.

\bibitem{FischlerRivoalEffective}
S.~Fischler and T.~Rivoal,
\emph{Effective algebraic independence of values of $E$-functions},
Math. Z. \textbf{305} (2023), article~48.
\href{https://doi.org/10.1007/s00209-023-03373-9}{doi:10.1007/s00209-023-03373-9}.

\bibitem{Kolchin}
E.~R.~Kolchin,
\emph{Differential Algebra and Algebraic Groups},
Pure and Applied Mathematics, vol.~54,
Academic Press, New York, 1973.

\bibitem{LongEIT}
C.~D.~Long and A.-A.~Le~Blanc,
\emph{Algebraic Exponential Integrals: Transcendence and Linear Relations},
companion preprint, August 2026.

\bibitem{LongLinearKZ}
C.~D.~Long and A.-A.~Le~Blanc,
\emph{A Linear Kontsevich--Zagier Theorem for Compact Exponential Periods on
the Affine Line},
companion preprint, August 2026.

\bibitem{LongPolynomialKZ}
C.~D.~Long and A.-A.~Le~Blanc,
\emph{A Polynomial Kontsevich--Zagier Theorem for Compact Exponential Periods
on the Affine Line},
companion preprint, August 2026.

\bibitem{RivoalRoques}
T.~Rivoal and J.~Roques,
\emph{On the algebraic dependence of $E$-functions},
Bull. Lond. Math. Soc. \textbf{48} (2016), no.~2, 271--279.
\href{https://doi.org/10.1112/blms/bdw003}{doi:10.1112/blms/bdw003}.

\bibitem{Salikhov}
V.~Kh.~Salikhov,
\emph{The algebraic independence of the values of $E$-functions satisfying
linear first-order differential equations},
Math. Notes Acad. Sci. USSR \textbf{13} (1973), 19--25.
\href{https://doi.org/10.1007/BF01093623}{doi:10.1007/BF01093623}.

\bibitem{vanDerPutSinger}
M.~van~der~Put and M.~F.~Singer,
\emph{Galois Theory of Linear Differential Equations},
Grundlehren der mathematischen Wissenschaften, vol.~328,
Springer, Berlin, 2003.
\href{https://doi.org/10.1007/978-3-642-55750-7}{doi:10.1007/978-3-642-55750-7}.

\end{thebibliography}
\end{document}
